Genealogical databases built by many contributors contain the same historical person many times.
In our setting, a database of roughly 150{,}000 people spread over 69 interconnected WebTrees family trees, the same Acadian ancestor is often recorded in several trees, each time with a slightly different name (\textit{Comeau} vs.\ \textit{Comeaux}, \textit{Marie} vs.\ \textit{Mary}), a dropped first given name, a \textit{dit} name (an inherited alias such as \textit{Jean Doucet dit Laverdure}), or approximate and missing dates.
Finding these duplicates matters to the genealogists who maintain the trees: merging them links branches that were researched independently, exposes contradictions between sources, and removes repeated work.
It is hard because Acadian naming is extremely concentrated (the ten most common surnames cover 39\% of all records, Section~\ref{sec:data}) and because the attributes that would separate two namesakes, such as birth dates and places, are often missing.

The method currently used on this database pools people who share an exact normalised name and counts concordant criteria (matching birth, marriage, spouse, parents, and so on).
It misses every duplicate whose name was written differently, and our own analysis of its output (Section~\ref{sec:experiment}) shows that most of its confident matches come from relatives (spouses, parents, children), not from the person's own record.
This suggests that the \emph{family structure} around a person carries much of the identifying signal, and motivates the question behind this project:
\emph{when does the family-graph structure around a person help, compared with that person's record alone?}

\textbf{Project type.} The project is both an \emph{improvement} project (we compare learned retrieval against the existing criteria-based method and attribute-only baselines) and an \emph{exploratory} one (we characterise when structure helps and when it does not).

\textbf{Objectives.} (i) Measure whether graph neural network embeddings of a person's family neighbourhood retrieve the true duplicate more reliably than attribute-only methods, and under which kinds of noise; (ii) measure whether giving a large language model (LLM) the family neighbourhood of a pair improves its match decision and the faithfulness of its justification.

\textbf{Input.} The person records of the database (names, sex, dated and placed events: birth, baptism, marriage, death, burial) and the typed family links between them (parent--child, spouse), exported from the WebTrees MySQL database through our Python library \texttt{genealogie\_acadienne} (Section~\ref{sec:data}).

\textbf{Expected outcome.} For retrieval, a ranked list of candidate duplicates for each query person; for verification, a decision (\textit{match}, \textit{not match}, \textit{insufficient information}) with cited evidence. Beyond the systems themselves, the expected outcome is an evidence-based answer to \emph{when} structure helps, per noise type.

\textbf{Success measurement.} Retrieval: Recall@$K$ ($K\in\{5,10\}$) and Mean Reciprocal Rank (MRR) on a benchmark with known duplicates. Verification: precision, recall and F1 of the \textit{match} decision on hand-labelled real pairs, plus the share of cited facts that actually exist in the given context (faithfulness).

\subsection*{Research questions}
Following the Milestone~1 feedback, we clarify that the two RQs are \emph{two separate experiments} asking the same underlying question in two settings, not two stages of one pipeline. Both contrast the same two conditions: the person's own record ($k=0$ hops) versus the record plus its family neighbourhood ($k=1$ or $2$ hops).

\textbf{RQ1 (retrieval):} \emph{Do relational graph convolutional network (R-GCN) embeddings of a person's $k$-hop family neighbourhood ($k=1,2$) rank the true duplicate higher than attribute-only methods, and under which kinds of noise?}
Retrieval is where the current method fails first: a duplicate that never enters the candidate pool can never be merged. The missing and noisy information is \emph{controlled}: we clone well-documented people into a dummy tree and apply one perturbation per tier (Section~\ref{sec:experiment}), so we know which kind of noise each result reflects.

\textbf{RQ2 (verification):} \emph{Given a candidate pair, does an LLM decide match / not match / insufficient more accurately, and justify its decision more faithfully, with the two people's 2-hop neighbourhoods ($k=2$) than with the two records alone ($k=0$)?}
Practitioners need a decision they can check, not only a score. RQ2 uses 150--200 hand-labelled real pairs, kept separate from the RQ1 clones; to connect the two experiments, the pairs are drawn from the top-ranked candidates of the RQ1 rankers on real data.
